% archon:covers AlgebraicJacobian/RiemannRoch/RationalCurveIso.lean
\chapter{The rational-curve isomorphism \(C \cong \mathbb P^1\) (RR.4)}
\label{chap:RiemannRoch_RationalCurveIso}
\uses{def:genus, lem:morphism_to_p1_from_global_sections, lem:degree_via_pole_divisor, lem:degree_one_morphism_iso, cor:lineBundleAtClosedPoint_exists_nonconstant_genusZero, thm:riemannRoch_genus_zero, def:codim1_cycles, def:divisor_degree}
% (read from references/hartshorne-algebraic-geometry.pdf, PDF page 314)
% + IV \S2, opening of "Hurwitz's Theorem" (degree of a finite morphism of
% curves is the function-field extension degree; read from PDF page 316)
% + I \S6 Corollary~6.12 (function-field \(\Leftrightarrow\) smooth-projective-curve
% equivalence of categories; read from PDF page 62). The verbatim quotes
% pasted in the declaration blocks below were copied character-by-character
% from the rendered page images of the scanned PDF (no text layer is present).
\section{Setup and motivation}
This chapter is the fourth and final sub-build chapter
\texttt{RR.1}--\texttt{RR.4} closing the project's headline \emph{RR bridge}
\[
  \mathtt{genusZero\_curve\_iso\_P1}
  \,\colon\,
  \text{a smooth proper geometrically irreducible curve \(C / \bar k\) with
  \(g(C) = 0\) is isomorphic to \(\mathbb P^1_{\bar k}\).}
\]
The chain assembled across the four sub-build chapters is:
\begin{itemize}
  \item \texttt{RR.1} (\texttt{chap:RiemannRoch_WeilDivisor}): the Weil-divisor
    group \(\Div(C)\), the degree map
    \(\deg \colon \Div(C) \to \mathbb Z\) (\texttt{def:divisor_degree}), and the
    degree-zero principal-divisor identity
    (\texttt{thm:principal_deg_zero}).
  \item \texttt{RR.2} (\texttt{chap:RiemannRoch_RRFormula}): the Riemann--Roch
    formula in genus zero \(\ell(D) = \deg(D) + 1\) for \(\deg(D) \geq 0\)
    (\texttt{thm:riemannRoch_genus_zero}).
  \item \texttt{RR.3} (sibling chapter, dispatched parallel iter-174):
    the line bundle \(\mathcal O_C([P])\), its global-sections identification
    \(\dim_{\bar k} H^0(C, \mathcal O_C([P])) = 2\), and the corollary
    \texttt{cor:nonconstant\_function\_genus\_zero} producing a non-constant
    rational function $f \in H^0(C, \mathcal O_C([P])) \setminus
    \bar k \cdot 1$ on a genus-$0$ curve with a marked $\bar k$-point.
  \item \texttt{RR.4} (\emph{this chapter}): the global classification ``a
    smooth proper geometrically irreducible curve \(C / \bar k\) with \(g(C) = 0\)
    and a \(\bar k\)-rational point is isomorphic to \(\mathbb P^1_{\bar k}\).''
\end{itemize}
This chapter consumes \texttt{RR.3} (the non-constant function) and produces
the Lean target \texttt{AlgebraicGeometry.genusZero\_curve\_iso\_P1} pinned at
\texttt{AlgebraicJacobian/AbelianVarietyRigidity.lean} line~290 and restated as
the interface block \texttt{chap:RiemannRoch_RationalCurveIso} of
\texttt{AbelianVarietyRigidity.tex}. The \texttt{RR.4}-closure of that target
also drops the \texttt{sorryAx} propagation tagged on
\texttt{rigidity\_genus0\_curve\_to\_grpScheme} via
\texttt{thm:rigidity_genus0_curve_to_AV}.
\paragraph{Scope of the \(\bar k\)-base classification.} The chapter restricts
to the case where the base field \(\bar k\) is algebraically closed. Over $\bar
k$ every closed point is $\bar k$-rational, so the ``$C$ has a $\bar
k$-rational point'' hypothesis of Hartshorne's Exercise~IV.1.3 is automatic
(the smooth proper geometrically irreducible curve \(C / \bar k\) is non-empty,
hence has closed points, all of which are \(\bar k\)-rational). The general
``\(k \to \bar k\) descent'' step ``genus-\(0\) curve over a non-closed \(k\) with
a \(k\)-rational point \(\cong \mathbb P^1_k\)'' is a separate sub-build hosted
in \texttt{Jacobian.tex} (the \texttt{genusZeroWitness.key} downstream
consumer) and is out of scope here.
\paragraph{Standing hypotheses.} Throughout this chapter, \(\bar k\) is an
algebraically closed field and \(C\) is a smooth proper geometrically
irreducible curve over \(\bar k\). The Lean signature
(\texttt{AlgebraicJacobian/RiemannRoch/RationalCurveIso.lean}) threads the
project's standing curve typeclass package:
\begin{quote}
  \texttt{\{kbar : Type u\} [Field kbar] [IsAlgClosed kbar]\\
   (C : Over (Spec (.of kbar))) [IsProper C.hom]\\
   \phantom{}\quad[SmoothOfRelativeDimension 1 C.hom]\\
   \phantom{}\quad[GeometricallyIrreducible C.hom]}
\end{quote}
matching the typeclass set used by the headline target
\texttt{AlgebraicGeometry.genusZero\_curve\_iso\_P1} in
\texttt{AlgebraicJacobian/AbelianVarietyRigidity.lean:290--296}. The genus
\(g(C)\) is the one defined in \texttt{def:genus} of \texttt{Genus.tex},
namely \(g(C) = \dim_{\bar k} H^1(C, \mathcal O_C)\) via the project's
\(\Module \bar k\)-flavoured sheaf cohomology
(\texttt{AlgebraicGeometry.genus}); the concrete projective line target
\(\mathbb P^1_{\bar k}\) is the project's
\texttt{AlgebraicGeometry.ProjectiveLineBar~kbar} of
\texttt{AlgebraicJacobian/Genus0BaseObjects.lean}.
\section{The morphism to \(\mathbb P^1\) from two global sections}
\begin{lemma}


\leanok
[Morphism to \(\mathbb P^1\) from a globally generating pair of sections]
  \label{lem:morphism_to_p1_from_global_sections}
  
  \lean{AlgebraicGeometry.Scheme.morphismToP1OfGlobalSections}
  Let \(X\) be a scheme over a field \(\bar k\), let \(\mathcal L\) be an invertible
  sheaf on \(X\), and let \(s_0, s_1 \in H^0(X, \mathcal L)\) be two global
  sections whose common zero locus is empty:
  \[
    \{x \in X : s_0(x) = 0 \text{ and } s_1(x) = 0\} \;=\; \emptyset.
  \]
  (Equivalently, \(s_0\) and \(s_1\) generate \(\mathcal L\) at every point.) Then
  there exists a unique morphism of \(\bar k\)-schemes
  \[
    \varphi_{(s_0, s_1)} \colon X \longrightarrow \mathbb P^1_{\bar k}
  \]
  whose pullback of \(\mathcal O_{\mathbb P^1}(1)\) is \(\mathcal L\) and whose
  pullbacks of the standard coordinate sections $x_0, x_1 \in
  H^0(\mathbb P^1, \mathcal O(1))$ are $s_0$ and $s_1$ respectively. On the
  basic-open chart $D_+(x_0) = \mathbb A^1_{\bar k} \subset
  \mathbb P^1_{\bar k}$ the morphism $\varphi_{(s_0, s_1)}$ is the map ``$x
  \mapsto s_1(x) / s_0(x)$''; on the chart $D_+(x_1)$ it is the map ``$x
  \mapsto s_0(x) / s_1(x)$''; the two definitions agree on the overlap
  \(D_+(x_0) \cap D_+(x_1) = \mathbb G_{m, \bar k}\) by the cocycle identity
  \((s_1 / s_0) \cdot (s_0 / s_1) = 1\).
\end{lemma}
\begin{proof}



\leanok
  This is the standard linear-system construction in the form of
  Hartshorne's Theorem~II.7.1 (``Morphisms to \(\mathbb P^n\)''), specialised to
  \(n = 1\). Mathlib (snapshot \texttt{b80f227}) ships the corresponding API
  under the name
  \texttt{AlgebraicGeometry.Proj.fromOfGlobalSections} in
  \texttt{Mathlib.AlgebraicGeometry.ProjectiveSpectrum.Basic}:
  \begin{quote}
  \small
  \texttt{def Proj.fromOfGlobalSections \{\(\sigma\) A\} [CommRing A]
    [SetLike \(\sigma\) A] [AddSubgroupClass \(\sigma\) A] (\(\mathcal A\) :
    \(\mathbb N\) → \(\sigma\)) [GradedRing \(\mathcal A\)] \{X : Scheme\}\\
    \phantom{}\quad(f : A →+* X.presheaf.obj \(\top\))\\
    \phantom{}\quad(hf : Ideal.map f (HomogeneousIdeal.irrelevant \(\mathcal A\)).toIdeal = \(\top\)) :\\
    \phantom{}\quad(X \(\longrightarrow\) Proj \(\mathcal A\))}
  \end{quote}
  Specialising to \(\mathcal A = \bar k[x_0, x_1]\) with the standard $\mathbb
  N$-grading by total degree, the inputs are a graded $\bar k$-algebra
  homomorphism \(f \colon \bar k[x_0, x_1] \to H^0(X, \mathcal O_X)\) and a
  proof that the irrelevant ideal \((x_0, x_1) \cdot \bar k[x_0, x_1]\) maps to
  the unit ideal of \(H^0(X, \mathcal O_X)\) (i.e.\ no common zeros). The
  resulting morphism lands in $\mathrm{Proj}\,\bar k[x_0, x_1] = \mathbb
  P^1_{\bar k}$ (Hartshorne~II, Example~7.1.3, the construction of $\mathbb
  P^n$ as a $\mathrm{Proj}$). When $\mathcal L$ is non-trivial the present
  lemma's signature is reduced to the Mathlib API by working on the open
  cover \(X = D(s_0) \cup D(s_1)\), where on each chart the sections become
  honest global functions (\(s_1 / s_0\) on \(D(s_0)\), \(s_0 / s_1\) on \(D(s_1)\),
  agreeing on the overlap). The chart-glue is performed by the
  project-bespoke wrapper
  \texttt{AlgebraicGeometry.Scheme.morphismToP1OfGlobalSections}, whose body
  in \texttt{AlgebraicJacobian/RiemannRoch/RationalCurveIso.lean} chains
  \texttt{Proj.fromOfGlobalSections} with the chart-restriction-and-glue
  pattern of \texttt{Scheme.Cover.glueMorphisms}.
  Uniqueness follows from the standard separation argument for morphisms to
  a separated scheme: two morphisms \(\varphi, \psi \colon X \to \mathbb P^1\)
  with \(\varphi^* x_i = \psi^* x_i\) for \(i = 0, 1\) agree on the basic-open
  charts \(D_+(x_0), D_+(x_1)\) (where they pull back the affine coordinates to
  the same rational function), and these charts cover \(\mathbb P^1\), so the
  two morphisms agree globally by gluing.
\end{proof}
\section{Degree of a morphism to \(\mathbb P^1\) via its pole divisor}
% NOTE (iter-190 plan-phase Pin 2 corrective; supersedes iter-182 narrative).
% `Scheme.Hom.poleDivisor \varphi` is the Weil divisor \(\varphi^*[\infty]\),
% defined as `positivePart (principal (\varphi^{\#} t_\infty))` in Lean.
% The supporting `Scheme.WeilDivisor.positivePart` def lives in
% `RiemannRoch_WeilDivisor.tex` (\ref{def:WeilDivisor_positivePart}); the
% Hartshorne~II.6.9 identity that powers the existential degree statement
% lives there as \ref{lem:degree_positivePart_principal_eq_finrank}.
\begin{lemma}


\leanok
[Degree of a non-constant morphism to \(\mathbb P^1\) via the preimage of \(\infty\)]
  \label{lem:degree_via_pole_divisor}
  \uses{def:WeilDivisor_positivePart,
        lem:degree_positivePart_principal_eq_finrank, def:principal_divisor,
        def:divisor_degree, thm:principal_deg_zero}
  \lean{AlgebraicGeometry.Scheme.morphism_degree_via_pole_divisor}
  % SOURCE: [Hartshorne], IV \S2, opening paragraph of "Hurwitz's Theorem"
  % (definition of the degree of a finite morphism of curves) (read from
  % references/hartshorne-algebraic-geometry.pdf, PDF page 316, doc p.~299).
  % SOURCE QUOTE: "Recall that the \emph{degree} of a finite morphism
  % \(f: X \to Y\) of curves is defined to be the degree \([K(X):K(Y)]\) of the
  % extension of function fields (II, \S 6)."
  \textit{Source: Hartshorne, IV.2, p.~299, opening of \S2.}
  Let \(\varphi \colon C \to \mathbb P^1_{\bar k}\) be a non-constant morphism
  from a smooth proper geometrically irreducible curve \(C / \bar k\). Then
  \(\varphi\) is finite, its function-field extension
  \(\bar k(\mathbb P^1) \hookrightarrow K(C)\) (the canonical
  \(\varphi\)-induced map, supplied as an
  \texttt{[Algebra (ProjectiveLineBar kbar).functionField C.functionField]}
  instance argument on the Lean side) is finite, and there exists a Weil
  divisor \(D \in \Div(C)\) such that
  \[
    D \;=\; \varphi^*[\infty]
    \quad\text{and}\quad
    \deg(D) \;=\; [\,K(C) : \bar k(\mathbb P^1)\,].
  \]
  Equivalently, the morphism \(\varphi\) satisfies
  \[
    [\,K(C) : \bar k(\mathbb P^1)\,] \;=\;
    \deg\!\left(\,\sum_{P \in \varphi^{-1}(\infty)}
                  \ord_P(\varphi^* t_\infty) \cdot [P]\,\right),
  \]
  where \(t_\infty \in K(\mathbb P^1)\) is a local parameter at \(\infty\)
  (equivalently \(t_\infty = 1/u\) for the standard affine coordinate \(u\) on
  \(\mathbb P^1 \setminus \{\infty\}\)); and for any closed point
  \(Q \in \mathbb P^1\) the same identification gives
  \(\deg(\varphi^*[Q]) = [K(C):\bar k(\mathbb P^1)]\).
  \emph{Convention on the pole divisor.} The Weil divisor
  \(\varphi^*[\infty]\) is exposed in Lean as
  \texttt{AlgebraicGeometry.Scheme.Hom.poleDivisor \(\varphi\)}, a
  noncomputable definition introduced (as a typed \texttt{sorry}) at the
  iter-182 plan-phase signature strengthening of this lemma. The body of
  \texttt{poleDivisor} is queued for iter-183+ closure via the standard
  affine-chart construction (\texttt{Ideal.sum\_ramification\_inertia}
  in Mathlib, applied at each closed point \(P \in \varphi^{-1}(\infty)\)
  to read off the local ramification multiplicity
  \(\ord_P(\varphi^* t_\infty)\) from the maximal ideal at \(\infty\) in
  \(\mathbb P^1_{\bar k}\)). On the Lean side, the existential output of
  \texttt{morphism\_degree\_via\_pole\_divisor} packages
  \(D = \texttt{Scheme.Hom.poleDivisor}\;\varphi\) together with
  \(\deg(D) = \operatorname{Module.finrank}_{K(\mathbb P^1)} K(C)\),
  mirroring the OCofP pattern in which a typed-sorry def
  (\texttt{def:lineBundleAtClosedPoint_toFunctionField}) is introduced
  upstream of its consumer lemma.
\end{lemma}
% NOTE (iter-190 plan-phase, supersedes iter-189 review NOTE):
% Pin 2 corrective COMMITTED — Option (a). `Scheme.Hom.poleDivisor` is
% being refactored to return the *positive part* of the principal divisor
% of `\varphi^{\#} t_\infty`, i.e. `(div(\varphi^{\#} t_\infty))_0
% = \varphi^*[\infty]`. The supporting Weil-divisor `positivePart`
% definition lives in `RiemannRoch_WeilDivisor.tex` as
% `def:WeilDivisor_positivePart`; the substantive Hartshorne~II.6.9
% identity that powers `lem:degree_via_pole_divisor` is now pinned as
% `lem:degree_positivePart_principal_eq_finrank` in `WeilDivisor.tex`.
% With these substrates in place, `lem:degree_via_pole_divisor` is
% mathematically correct as stated: the body's existential witness `D`
% is `Hom.poleDivisor \varphi = positivePart (principal (\varphi^{\#}
% t_\infty))`, and the degree identity is a one-line corollary of
% `lem:degree_positivePart_principal_eq_finrank`. iter-190 prover lane
% (Lane I) implements the refactor:
%   1. New def `Scheme.WeilDivisor.positivePart` in `WeilDivisor.lean`
%      (~10-15 LOC; lattice `D \sqcup 0`).
%   2. New lemma
%      `Scheme.WeilDivisor.degree_positivePart_principal_eq_finrank` as a
%      named typed-sorry pin in `WeilDivisor.lean` (~5 LOC signature; body
%      ~50-80 LOC via affine-chart `Ideal.sum_ramification_inertia` route
%      owed iter-191+).
%   3. Refactor `Scheme.Hom.poleDivisor` body in `RationalCurveIso.lean`
%      from `principal (algebraMap _ _ t.val) halg` to
%      `positivePart (principal (algebraMap _ _ t.val) halg)` (~5 LOC).
%   4. Close `Hom.poleDivisor_degree_eq_finrank` body via a 1-line
%      rewrite consuming (2) (~3 LOC).
% Net iter-190 Lane I: 1 axiom-clean closure (Pin 2 body, modulo named
% typed sorry on (2)); 1 NEW named typed-sorry pin
% (`degree_positivePart_principal_eq_finrank`) which is itself an
% iter-191+ Hartshorne II.6.9 substrate target. See iter-190 sidecar
% `iter/iter-190/plan.md` for the full rationale.
\begin{proof}
  \uses{def:WeilDivisor_positivePart, lem:degree_positivePart_principal_eq_finrank, def:principal_divisor, def:divisor_degree, thm:principal_deg_zero}
  \leanok
  A non-constant morphism between smooth proper curves over \(\bar k\) is
  automatically finite. Indeed, the image of \(\varphi\) is a closed
  irreducible subset of \(\mathbb P^1_{\bar k}\) that is not a point (else
  \(\varphi\) would be constant) and not empty, so it is all of $\mathbb
  P^1_{\bar k}$ (the only proper irreducible closed subset of $\mathbb
  P^1_{\bar k}$ is a point); hence $\varphi$ is surjective. A surjective
  morphism between smooth proper irreducible curves of equal dimension is
  finite (proper + quasi-finite, where quasi-finiteness comes from the
  fact that \(\varphi\) has zero-dimensional fibres: a positive-dimensional
  fibre would force a one-dimensional irreducible component in the fibre,
  contradicting irreducibility of \(\mathbb P^1_{\bar k}\) and \(C\) at
  generic points). Hence the function-field extension
  \(\bar k(\mathbb P^1) = \varphi^\# K(\mathbb P^1) \hookrightarrow K(C)\)
  is finite.
  The identification \(\deg(\varphi^*[Q]) = [K(C) : \bar k(\mathbb P^1)]\) for
  any closed point \(Q \in \mathbb P^1\) is the multiplicativity of the
  degree under finite pullback (Hartshorne~II.6.9): for a finite morphism
  \(\varphi \colon C \to C'\) of smooth proper curves and any divisor $D \in
  \Div(C')$, $\deg(\varphi^* D) = \deg(\varphi) \cdot \deg(D)$. Specialising
  to \(D = [Q]\) (a divisor of degree \(1\) since \(\kappa(Q) = \bar k\)) gives
  \(\deg(\varphi^*[Q]) = \deg(\varphi)\). Equivalently, the degree of
  \(\varphi\) equals the number of preimages of any closed point of $\mathbb
  P^1$, counted with the ramification multiplicities $\ord_P(\varphi^*
  t_\infty)$ at the preimage points $P \in \varphi^{-1}(\infty)$.
  \emph{Sub-build note.} The lemma threads two auxiliary inputs:
  ``non-constant morphism between smooth proper curves is finite'' (the
  surjective-and-proper-and-quasi-finite argument above; Hartshorne~II.6.8)
  and ``multiplicativity of degree under finite pullback'' (Hartshorne~II.6.9
  and the multiplicativity of \(\deg\) under composition of \(\div\) and
  \(\deg\); the closely related ``\(\mathrm{div}(f)\) has degree zero'' specialisation
  is \texttt{thm:principal_deg_zero} of \texttt{RR.1}). Both are separable
  sub-lemmas inside \texttt{RR.4}'s Lean file
  (\texttt{AlgebraicJacobian/RiemannRoch/RationalCurveIso.lean}) and are
  slotted for closure in a follow-up iteration after the
  \texttt{morphism\_degree\_via\_pole\_divisor} skeleton lands.
\end{proof}
\subsection{The pole divisor \(\varphi^*[\infty]\) and its degree (Lean substrate)}
% The four blocks of this subsection give 1-to-1 blueprint coverage for the
% Lean helper declarations underlying `lem:degree_via_pole_divisor` above
% (the local parameter at \(\infty\), its uniformiser property, the pole
% divisor itself, and the function-field-degree identity). They are
% project-internal plumbing; the external grounding is the Hartshorne IV.2 /
% II.6.9 material already cited on `lem:degree_via_pole_divisor`, so no
% further verbatim citation is attached here.
\begin{definition}[Local parameter at \(\infty \in \mathbb P^1\)]
  \label{def:localParameterAtInfty}
  \lean{AlgebraicGeometry.Scheme.localParameterAtInfty}
  \uses{def:functionFieldIso}
  Let \(\bar k\) be a field and present the projective line as
  \(\mathbb P^1_{\bar k} = \operatorname{Proj} \bar k[X_0, X_1]\). The
  \emph{local parameter at \(\infty\)} is the rational function
  \[
    t_\infty \;:=\; X_0 / X_1 \;\in\; K(\mathbb P^1_{\bar k}),
  \]
  the degree-\(0\) ratio coordinate of the homogeneous localisation
  \(\bigl(\bar k[X_0, X_1]_{(X_1)}\bigr)_0 = \mathcal O\!\bigl(D_+(X_1)\bigr)\),
  transported into the function field \(K(\mathbb P^1)\) along the germ map
  on the affine chart \(D_+(X_1)\). On the Lean side it is packaged as a
  \emph{nonzero} element of \(K(\mathbb P^1)\): the witness \(t_\infty \neq 0\)
  comes from the injectivity of the chart isomorphism
  \(\mathcal O\!\bigl(D_+(X_1)\bigr) \cong \bar k[X_0/X_1]\) (a
  \(\operatorname{Proj}\)-chart presentation) composed with the injective
  germ-to-function-field map of the integral scheme \(\mathbb P^1_{\bar k}\),
  applied to the nonzero ratio element \(X_0 / X_1\).
\end{definition}
\begin{proof}
  Proved directly in Lean.
\end{proof}
\begin{lemma}[\(t_\infty\) is a uniformiser at \([\infty]\)]
  \label{lem:localParameterAtInfty_uniformiser_witness}
  \uses{def:localParameterAtInfty}
  \lean{AlgebraicGeometry.Scheme.localParameterAtInfty_uniformiser_witness}
  There is a unique prime divisor \(Y_0 \in \Div(\mathbb P^1_{\bar k})\) — the
  point at infinity \([\infty]\) — at which the local parameter
  \(t_\infty = X_0 / X_1\) of \cref{def:localParameterAtInfty} has positive
  order, and there \(\ord_{Y_0}(t_\infty) = 1\). Equivalently, the zero divisor
  \((t_\infty)_0 = [\infty]\) of \(t_\infty\) on \(\mathbb P^1_{\bar k}\) has
  degree \(1\) and is supported at the single prime divisor \([\infty]\).
\end{lemma}
\begin{proof}
  \uses{def:localParameterAtInfty}
  \emph{Honest sketch (the Lean body is an open hole).} The witness \(Y_0\) is
  the closed point at infinity, image under \(\operatorname{Proj}.\mathrm{aw\!ay\iota}\)
  at \(X_1\) of the maximal ideal generated by \(X_0/X_1\) in the affine chart
  \(\operatorname{Spec} \bar k[X_0/X_1]\). On that chart the stalk at \(Y_0\) is
  the discrete valuation ring \(\bar k[X_0/X_1]_{(X_0/X_1)}\) with uniformiser
  \(X_0/X_1 = t_\infty\), so \(\ord_{Y_0}(t_\infty) = 1\); and on the
  complementary chart \(D_+(X_0)\) the function \(t_\infty\) is a unit, whence
  no other prime divisor carries positive order. Formalising this requires the
  not-yet-available stalk-of-\(\operatorname{Proj}\)-chart-is-a-DVR bridge and
  the order-via-uniformiser identity (the same codimension-one regularity
  substrate that \(\mathbb P^1_{\bar k}\)'s
  \texttt{IsRegularInCodimensionOne} instance is gated on); pending that, the
  Lean proof is left as an explicit hole.
\end{proof}
\begin{definition}[Pole divisor of \(\varphi \colon C \to \mathbb P^1\)]
  \label{def:hom_poleDivisor}
  \uses{def:localParameterAtInfty, def:WeilDivisor_positivePart, def:principal_divisor}
  \lean{AlgebraicGeometry.Scheme.Hom.poleDivisor}
  Let \(C / \bar k\) be a smooth proper geometrically irreducible curve, regular
  in codimension one, and let \(\varphi \colon C \to \mathbb P^1_{\bar k}\) be a
  morphism, with \(K(C)\) carrying the canonical \(\varphi\)-induced
  \(K(\mathbb P^1)\)-algebra structure. The \emph{pole divisor} of \(\varphi\)
  is the Weil divisor
  \[
    \varphi^*[\infty]
    \;:=\;
    \bigl(\operatorname{div}(\varphi^\# t_\infty)\bigr)_0
    \;=\;
    \operatorname{positivePart}\!\bigl(\operatorname{div}(\varphi^\# t_\infty)\bigr)
    \;\in\; \Div(C),
  \]
  the positive part of the principal divisor on \(C\) of the pulled-back local
  parameter \(\varphi^\# t_\infty = \operatorname{algebraMap}_{K(\mathbb P^1) \to K(C)}(t_\infty)\),
  where \(t_\infty\) is the local parameter of \cref{def:localParameterAtInfty}.
  The pullback is nonzero because the algebra map \(K(\mathbb P^1) \to K(C)\) is
  injective (a ring homomorphism out of a field), so \(t_\infty \neq 0\) gives
  \(\varphi^\# t_\infty \neq 0\) and the principal divisor is well defined. On a
  complete nonsingular curve the zeros of \(\varphi^\# t_\infty\) are exactly the
  preimages of \(\infty\), so the positive part recovers \(\varphi^*[\infty]\) on
  the nose.
\end{definition}
\begin{proof}
  Proved directly in Lean.
\end{proof}
\begin{theorem}[Degree of the pole divisor equals the function-field degree]
  \label{thm:poleDivisor_degree_eq_finrank}
  \uses{def:hom_poleDivisor, lem:localParameterAtInfty_uniformiser_witness,
        lem:degree_positivePart_principal_eq_finrank, def:divisor_degree}
  \lean{AlgebraicGeometry.Scheme.Hom.poleDivisor_degree_eq_finrank}
  For a non-constant morphism \(\varphi \colon C \to \mathbb P^1_{\bar k}\) of
  smooth proper geometrically irreducible curves over \(\bar k\) (with \(K(C)\)
  finite over \(K(\mathbb P^1)\) via the \(\varphi\)-induced algebra structure),
  the degree of the pole divisor equals the function-field extension degree:
  \[
    \deg\!\bigl(\varphi^*[\infty]\bigr)
    \;=\;
    [\,K(C) : K(\mathbb P^1)\,]
    \;=\;
    \operatorname{finrank}_{K(\mathbb P^1)} K(C).
  \]
  This is the precise form of the degree identity packaged existentially by
  \cref{lem:degree_via_pole_divisor}.
\end{theorem}
\begin{proof}
  \uses{def:hom_poleDivisor, lem:localParameterAtInfty_uniformiser_witness,
        lem:degree_positivePart_principal_eq_finrank}
  Proved directly in Lean: unfolding \(\varphi^*[\infty] =
  \operatorname{positivePart}(\operatorname{div}(\varphi^\# t_\infty))\)
  reduces the claim to the Hartshorne~II.6.9 degree identity
  \cref{lem:degree_positivePart_principal_eq_finrank} of \texttt{RR.1},
  specialised at the local parameter \(t = t_\infty\), whose uniformiser
  hypothesis is supplied by \cref{lem:localParameterAtInfty_uniformiser_witness}.
\end{proof}
\section{A degree-\(1\) morphism between smooth proper curves is an isomorphism}
% NOTE (iter-181 Lane I signature refinement; see analogies/ratcurveiso-pin3.md
% Decision 2, verdict DIVERGE_INTENTIONALLY). The iter-177 file-skeleton
% placeholder hypothesis "Nonempty (C'.functionField ≃+* C.functionField)"
% (an abstract function-field iso, unrelated to \(\varphi\)) was strictly
% weaker than the birational-extension body needs and was refined in
% iter-181 to the explicit function-field-extension-degree formulation
% [K(C):K(C')] = 1, i.e. Module.finrank C'.functionField C.functionField
% = 1, with C.functionField a C'.functionField-algebra via the canonical
% \varphi-induced map (Scheme.Hom.stalkMap at the generic point composed
% with IsFractionRing.lift). The proof-prose body below (the Hartshorne
% I.6.12 chain + scheme-theoretic alternative argument) is unchanged ---
% only the input shape changes.
% NOTE (iter-194 reviewer; refreshed iter-196): The Lean
% `iso_of_degree_one` body inline Step 1 was refactored iter-194 from
% 16 LOC of `Subalgebra.bot_eq_top_of_finrank_eq_one +
% Algebra.surjective_algebraMap_iff + RingHom.injective +
% RingEquiv.ofBijective` to 4 LOC consuming two new axiom-clean substrate
% helpers (`algebraMap_bijective_of_finrank_one` +
% `phi_left_functionField_algEquiv_of_finrank_one`).  Helper (a)
% `phi_left_locallyQuasiFinite_of_finrank_one` was reformulated iter-196
% via Mathlib's `LocallyQuasiFinite.of_finite_preimage_singleton`
% (`QuasiFinite.lean:295`) with in-scope derivation of `IsProper φ.left`
% + `LocallyOfFiniteType φ.left`; the residual gap is now the concrete
% topological claim `(φ.left ⁻¹' {x}).Finite` for arbitrary
% `x : C'.left`. The generic-point branch is closeable via
% `genericPoint_eq` functoriality + the algebra-iso substrate; the
% closed-point branch remains a Mathlib gap (\textit{smooth-dim-1}
% $\Rightarrow$ 0-dim fibre; iter-197+ candidate). Helper (d)
% `phi_left_fromNormalization_isIso_of_smoothProper_finrank_one` is
% gated on project-side `IsNormalScheme` substrate.
\begin{lemma}


\leanok
[Degree-\(1\) morphism between smooth proper curves is an isomorphism]
  \label{lem:degree_one_morphism_iso}
  \uses{lem:degree_via_pole_divisor, def:functionFieldIso, lem:phi_left_toNormalization_isIso_of_isIntegralHom, lem:phi_left_fromNormalization_isIso_of_smoothProper_finrank_one}
  \lean{AlgebraicGeometry.Scheme.iso_of_degree_one}
  % SOURCE: [Hartshorne], I \S6 Corollary~6.12 (function-field-to-curve
  % equivalence of categories) (read from
  % references/hartshorne-algebraic-geometry.pdf, PDF page 62, doc p.~45)
  % + [Hartshorne], IV \S2 opening paragraph (degree of a finite morphism
  % of curves is the function-field extension degree) (read from
  % references/hartshorne-algebraic-geometry.pdf, PDF page 316, doc p.~299).
  % SOURCE QUOTE: "Corollary 6.12. \emph{The following three categories are
  % equivalent:}
  % (i) \emph{nonsingular projective curves, and dominant morphisms;}
  % (ii) \emph{quasi-projective curves, and dominant rational maps;}
  % (iii) \emph{function fields of dimension 1 over \(k\), and \(k\)-homomorphisms.}"
  % SOURCE QUOTE: "Recall that the \emph{degree} of a finite morphism
  % \(f: X \to Y\) of curves is defined to be the degree \([K(X):K(Y)]\) of the
  % extension of function fields (II, \S 6)."
  \textit{Source: Hartshorne, I.6 Corollary~6.12 (p.~45) + IV.2 opening
  (p.~299).}
  Let \(C, C'\) be smooth proper geometrically irreducible curves over an
  algebraically closed field \(\bar k\), and let \(\varphi \colon C \to C'\) be
  a non-constant morphism whose induced function-field extension is of
  degree one,
  \[
    [\,K(C) : K(C')\,] \;=\; 1,
  \]
  formally
  \(\operatorname{Module.finrank}_{K(C')} K(C) = 1\) in the Lean signature.
  Then \(\varphi\) is an isomorphism of \(\bar k\)-schemes.
  \emph{Convention on the function-field map.} At each call-site of this
  lemma, the \(K(C')\)-algebra structure on \(K(C)\) used to interpret the
  degree \([K(C) : K(C')]\) is the canonical \(\varphi\)-induced
  function-field map, i.e.\ the composite
  \[
    K(C') \;\xrightarrow{\;\text{germ}\;}\;
    \mathcal O_{C', \varphi(\xi_C)}
    \;\xrightarrow{\;\varphi^\#\;}\;
    \mathcal O_{C, \xi_C}
    \;\xrightarrow{\;\text{frac.}\,\text{lift}\;}\; K(C)
  \]
  of the germ at the image of the generic point \(\xi_C \in C\) with the
  stalk-map of \(\varphi\) at \(\xi_C\) and the fraction-field lift at
  \(\xi_C\) (per \texttt{analogies/ratcurveiso-pin3.md}, Decision~2). On
  the Lean side this canonical instance is carried as a
  \texttt{[Algebra C'.functionField C.functionField]} instance argument
  to \texttt{AlgebraicGeometry.Scheme.iso\_of\_degree\_one}; consumers
  supply it explicitly at each call-site.
\end{lemma}
% NOTE (iter-189 review): Pin 3 Step 1 (function-field iso K(C') ≃+* K(C)
% from `Module.finrank = 1`) closed axiom-clean inline iter-189 (~10 LOC via
% `Subalgebra.bot_eq_top_of_finrank_eq_one` + `Algebra.surjective_algebraMap_iff`
% + `RingHom.injective` + `RingEquiv.ofBijective`). Step 2 (scheme-level lift
% via `Scheme.Hom.toNormalization`) remains a typed sorry in the body, with 4
% sub-obligations (a)-(d) documented inline at the sorry. Iter-190+ closure
% budget ~80-150 LOC per `analogies/ratcurveiso-pin3.md`.
\begin{proof}



\leanok
  The function-field hypothesis \([K(C) : K(C')] = 1\) states that the
  canonical \(\varphi\)-induced map \(\varphi^\# \colon K(C') \to K(C)\) is a
  finite extension of degree one; since \(K(C')\) is a field and the
  extension has dimension one as a \(K(C')\)-vector space, \(\varphi^\#\) is
  an isomorphism of function fields \(K(C') \cong K(C)\). (Equivalently, by
  the IV.2-opening identification \(\deg(\varphi) = [K(C):K(C')]\), this is
  the statement \(\deg(\varphi) = 1\); the morphism \(\varphi\) is finite in
  this case by the standard argument that a non-constant morphism of
  smooth proper curves is finite, identical to the
  \texttt{lem:degree_via_pole_divisor} preamble.)
  By Hartshorne's Corollary~I.6.12, the assignment \(C \mapsto \bar k(C)\) is
  an equivalence of categories between smooth projective curves over $\bar
  k$ (with dominant morphisms) and function fields of transcendence degree
  one over \(\bar k\) (with \(\bar k\)-homomorphisms). Under this equivalence,
  the isomorphism \(\bar k(C') \cong \bar k(C)\) produced by the previous
  paragraph lifts to a unique isomorphism of smooth proper curves $C \to
  C'$, and this isomorphism agrees with $\varphi$ on the generic point (by
  construction: it is the inverse-image of \(\varphi^\#\) under the
  equivalence). The two morphisms \(C \to C'\) that agree on the generic
  point of the integral scheme \(C\) agree everywhere (a morphism out of an
  integral scheme to a separated scheme is determined by its restriction
  to the generic point, Hartshorne~II Ex.~4.2). Hence \(\varphi\) is itself
  the lift, and is therefore an isomorphism.
  \emph{Alternative direct argument.} Without invoking the
  equivalence-of-categories framing, the same conclusion is obtained
  scheme-theoretically: a finite degree-\(1\) morphism $\varphi \colon C \to
  C'$ between integral normal noetherian schemes is an isomorphism.
  Indeed, the direct image \(\varphi_* \mathcal O_C\) is a coherent
  \(\mathcal O_{C'}\)-module that is generically of rank $\deg(\varphi) =
  1$, and an inclusion $\mathcal O_{C'} \hookrightarrow \varphi_*
  \mathcal O_C$ of coherent rank-$1$ sheaves whose generic ranks agree
  is an isomorphism over the locus where the quotient
  \(\varphi_* \mathcal O_C / \mathcal O_{C'}\) has trivial stalk. The
  quotient is a torsion coherent \(\mathcal O_{C'}\)-module of generic rank
  \(\deg(\varphi) - 1 = 0\), hence is itself torsion-free (a
  finitely-generated torsion module over a Dedekind domain of generic
  rank zero is zero), so vanishes. Therefore $\varphi_* \mathcal O_C =
  \mathcal O_{C'}$ and $\varphi$ is the structure morphism of an
  isomorphism. Cf.\ Stacks tag~\textbf{0AVX} for the categorical version.
  \emph{Sub-build note.} The Lean signature
  \texttt{AlgebraicGeometry.Scheme.iso\_of\_degree\_one} consumes a
  non-constant morphism \(\varphi \colon C \to C'\) together with a
  \texttt{[Algebra C'.functionField C.functionField]} instance argument
  (the canonical \(\varphi\)-induced function-field map) and the
  function-field-extension-degree-one hypothesis
  \texttt{Module.finrank C'.functionField C.functionField = 1}, and
  produces an isomorphism \(C \cong C'\) in the slice category
  \(\mathtt{Over (Spec (.of\ kbar))}\). The proof body threads the Mathlib
  API \texttt{IsIntegralHom.iso\_of\_isUnit} (or the equivalent
  \texttt{IsAffineHom + Finite + RankOne} chain) and, where needed, the
  project-bespoke ``finite-morphism-of-degree-one is an isomorphism''
  helper. The Hartshorne~I.6.12 equivalence-of-categories framing is the
  intended high-level proof structure; the scheme-theoretic argument
  above is the spelled-out form for the Lean body.
\end{proof}
\subsection{The normalization factorisation helpers (Lean substrate)}
% The five blocks of this subsection give 1-to-1 blueprint coverage for the
% Lean helper declarations into which the body of `lem:degree_one_morphism_iso`
% (`iso_of_degree_one`) decomposes: the bijectivity of the function-field
% algebra map under finrank one, the induced function-field iso, local
% quasi-finiteness of the underlying morphism, and the two normalization
% factors. They are project-internal plumbing whose external grounding is the
% Hartshorne I.6.12 / Stacks 0AVX material already cited on
% `lem:degree_one_morphism_iso`.
\begin{lemma}[Degree-one algebra map over a field is bijective]
  \label{lem:algebraMap_bijective_of_finrank_one}
  \lean{AlgebraicGeometry.Scheme.algebraMap_bijective_of_finrank_one}
  Let \(K\) be a field and \(L\) a \(K\)-algebra with
  \(\operatorname{finrank}_K L = 1\). Then the structure map
  \(\operatorname{algebraMap} \colon K \to L\) is bijective.
\end{lemma}
\begin{proof}
  Proved directly in Lean. A subalgebra of dimension one over a field is the
  whole algebra, i.e. \(\bot = \top\) as \(K\)-subalgebras of \(L\); over a field
  this is equivalent to the structure map being bijective.
\end{proof}
\begin{definition}[Function-field isomorphism from a degree-one extension]
  \label{def:phi_left_functionField_algEquiv_of_finrank_one}
  \uses{lem:algebraMap_bijective_of_finrank_one}
  \lean{AlgebraicGeometry.Scheme.phi_left_functionField_algEquiv_of_finrank_one}
  For integral schemes \(C, C'\) over \(\bar k\) with \(K(C)\) a
  \(K(C')\)-algebra of dimension one
  (\(\operatorname{finrank}_{K(C')} K(C) = 1\)), the structure map is the ring
  isomorphism
  \[
    K(C') \;\xrightarrow{\ \sim\ }\; K(C)
  \]
  obtained from the bijectivity of \cref{lem:algebraMap_bijective_of_finrank_one}.
\end{definition}
\begin{proof}
  Proved directly in Lean.
\end{proof}
\begin{lemma}[Local quasi-finiteness of a degree-one morphism]
  \label{lem:phi_left_locallyQuasiFinite_of_finrank_one}
  \uses{def:phi_left_functionField_algEquiv_of_finrank_one}
  \lean{AlgebraicGeometry.Scheme.phi_left_locallyQuasiFinite_of_finrank_one}
  Let \(C, C'\) be smooth proper geometrically irreducible curves over an
  algebraically closed field \(\bar k\) and \(\varphi \colon C \to C'\) a
  morphism whose induced function-field extension has degree one
  (\(\operatorname{finrank}_{K(C')} K(C) = 1\)). Then the underlying scheme
  morphism \(\varphi\) is locally quasi-finite. Combined with properness this
  upgrades to finiteness (and hence to an integral, affine morphism).
\end{lemma}
\begin{proof}
  \uses{def:phi_left_functionField_algEquiv_of_finrank_one}
  \emph{Honest sketch (the Lean body is an open hole).} By properness of
  \(\varphi\) (read off from the triangle over \(\operatorname{Spec} \bar k\))
  it suffices, via the reduction to fibre-set finiteness, to show every fibre
  \(\varphi^{-1}(x)\) is finite. At the generic point of \(C'\) the fibre is
  \(\operatorname{Spec} K(C) \to \operatorname{Spec} K(C')\), a single point
  because the function-field map is the isomorphism of
  \cref{def:phi_left_functionField_algEquiv_of_finrank_one}. At a closed point
  the fibre is zero-dimensional because both curves are smooth of relative
  dimension one over \(\bar k\). The closed-point case rests on the
  ``smooth of relative dimension one \(\Rightarrow\) zero-dimensional fibre''
  statement, which is not yet available in the project; pending it, the Lean
  proof is left as an explicit hole.
\end{proof}
\begin{lemma}[Normalization-target factor is an isomorphism]
  \label{lem:phi_left_toNormalization_isIso_of_isIntegralHom}
  \uses{lem:phi_left_locallyQuasiFinite_of_finrank_one}
  \lean{AlgebraicGeometry.Scheme.phi_left_toNormalization_isIso_of_isIntegralHom}
  For a quasi-compact, quasi-separated, integral morphism
  \(\varphi \colon C \to C'\) of schemes, the dominant factor of the
  normalization factorisation,
  \(\varphi.\mathrm{toNormalization} \colon C \to \widetilde{C'}_\varphi\), is an
  isomorphism. (The integrality hypothesis is the output of
  \cref{lem:phi_left_locallyQuasiFinite_of_finrank_one} together with
  properness.)
\end{lemma}
\begin{proof}
  \uses{lem:phi_left_locallyQuasiFinite_of_finrank_one}
  Proved directly in Lean: this is the Mathlib instance attaching
  \texttt{IsIso} to \(\varphi.\mathrm{toNormalization}\) for an integral
  quasi-compact quasi-separated morphism (the dominant half of the
  Hartshorne~I.6.12 / Stacks~\textbf{0AVX} factorisation).
\end{proof}
\begin{lemma}[Normalization-source factor is an isomorphism over a smooth target]
  \label{lem:phi_left_fromNormalization_isIso_of_smoothProper_finrank_one}
  \uses{def:phi_left_functionField_algEquiv_of_finrank_one}
  \lean{AlgebraicGeometry.Scheme.phi_left_fromNormalization_isIso_of_smoothProper_finrank_one}
  Let \(C'\) be a smooth proper geometrically irreducible curve over \(\bar k\),
  let \(C\) be integral, and let \(\varphi \colon C \to C'\) be a quasi-compact
  quasi-separated morphism whose induced function-field extension has degree one
  (\(\operatorname{finrank}_{K(C')} K(C) = 1\)). Then the integral factor of the
  normalization factorisation,
  \(\varphi.\mathrm{fromNormalization} \colon \widetilde{C'}_\varphi \to C'\), is
  an isomorphism.
\end{lemma}
\begin{proof}
  \uses{def:phi_left_functionField_algEquiv_of_finrank_one}
  \emph{Honest sketch (the Lean body is an open hole).} A smooth proper curve
  over \(\bar k\) is regular, hence normal, so its affine sections are integrally
  closed (Dedekind) domains. Under the degree-one hypothesis the function-field
  iso of \cref{def:phi_left_functionField_algEquiv_of_finrank_one} identifies
  \(K(C)\) with \(K(C')\), so the integral closure of each \(\Gamma(C', U)\)
  inside \(\Gamma(C, \varphi^{-1}U)\) collapses to \(\Gamma(C', U)\) itself; thus
  \(\varphi.\mathrm{fromNormalization}\) is an isomorphism. Formalising this needs
  a scheme-level normality substrate (``smooth curve over a field has Dedekind
  sections / is a normal scheme'') not yet present in the project; pending it,
  the Lean proof is left as an explicit hole.
\end{proof}
\section{Headline: genus-\(0\) curve over \(\bar k\) is isomorphic to \(\mathbb P^1\)}
\begin{theorem}

\leanok
[Genus-\(0\) curve over \(\bar k\) is isomorphic to \(\mathbb P^1\)]
  \label{thm:genus_zero_curve_iso_p1}
  \uses{thm:riemannRoch_genus_zero, cor:lineBundleAtClosedPoint_exists_nonconstant_genusZero,
        lem:morphism_to_p1_from_global_sections, lem:degree_via_pole_divisor,
        lem:degree_one_morphism_iso, def:genus}
  \lean{AlgebraicGeometry.genusZero_curve_iso_P1}
  % already know that \(p_a(\mathbf P^1) = 0\) (I, Ex. 7.2), so conversely
  % suppose given a curve \(X\) of genus 0. Let \(P, Q\) be two distinct
  % points of \(X\) and apply Riemann-Roch to the divisor \(D = P - Q\).
  % Since \(\deg(K - D) = -2\), using (1.3.3) above, we have \(l(K - D) = 0\),
  % and so we find \(l(D) = 1\). But \(D\) is a divisor of degree 0, so by
  % (1.2) we have \(D \sim 0\), in other words \(P \sim Q\). But this implies
  % \(X\) is rational (II, 6.10.1)."
  \textit{Source: Hartshorne, IV.1, Example~1.3.5, p.~297.}
  Let \(\bar k\) be an algebraically closed field, and let \(C\) be a smooth
  proper geometrically irreducible curve over \(\bar k\) with \(g(C) = 0\).
  Then there exists an isomorphism of \(\bar k\)-schemes
  \[
    \varphi \colon C \stackrel{\sim}{\longrightarrow} \mathbb P^1_{\bar k}.
  \]
  Concretely, the Lean signature
  \texttt{AlgebraicGeometry.genusZero\_curve\_iso\_P1} (the existing pinned
  target at \texttt{AlgebraicJacobian/AbelianVarietyRigidity.lean:290}) is
  closed by the proof below; the existing interface block
  \texttt{chap:RiemannRoch_RationalCurveIso} of
  \texttt{chap:avr_for_rr} is its consumer-side restatement,
  pinning the same Lean declaration.
\end{theorem}
% SOURCE QUOTE PROOF: see the verbatim Example IV.1.3.5 text in the
% statement block above (Hartshorne states the classification and its
% proof inside the single Example; the proof is the second half beginning
% "We already know that \(p_a(\mathbf P^1) = 0\) \ldots"). The present
% proof body restates Hartshorne's argument in the project's notation
% (linear-system construction \(C \to \mathbb P^1\) via $H^0(C, \mathcal O_C
% ([P]))\(, degree-\)1$ identification, isomorphism).
\begin{proof}
  Over the algebraically closed field \(\bar k\) the smooth proper
  geometrically irreducible curve \(C\) has closed points, all of which are
  \(\bar k\)-rational (the residue field of a closed point of a curve over
  \(\bar k\) is a finite extension of \(\bar k\), hence is \(\bar k\) itself by
  algebraic closedness). Fix one such \(\bar k\)-rational point \(P \in C\);
  this is the marked point used to build the linear system below.
  \emph{Step 1: non-constant function \(f \in H^0(C, \mathcal O_C([P]))\).}
  By Hartshorne's Riemann--Roch in genus zero
  (\texttt{thm:riemannRoch_genus_zero} of \texttt{RR.2}) applied to the
  effective divisor \([P]\) of degree \(1\), we obtain
  \(\ell([P]) = \deg([P]) + 1 = 2\). The constant section
  \(1 \in \bar k \cdot 1 \subset H^0(C, \mathcal O_C([P]))\) accounts for a
  one-dimensional \(\bar k\)-subspace; the two-dimensional quotient supplies
  a non-constant section $f \in H^0(C, \mathcal O_C([P])) \setminus
  \bar k \cdot 1$. This is the corollary
  \texttt{cor:nonconstant\_function\_genus\_zero} of \texttt{RR.3} (sibling
  chapter, \texttt{RiemannRoch\_OCofP.tex}).
  \emph{Step 2: morphism \(C \to \mathbb P^1\) via the linear system \((1, f)\).}
  The pair \((1, f) \in H^0(C, \mathcal O_C([P]))^2\) generates the
  invertible sheaf \(\mathcal O_C([P])\) at every point of \(C\): at points
  \(Q \in C \setminus \{P\}\) the section \(1\) does not vanish in
  \(\mathcal O_C([P])_Q \cong \mathcal O_{C, Q}\), and at the marked point
  \(P\) the section \(f\) does not vanish in \(\mathcal O_C([P])_P\) (because \(f\)
  has a pole of order exactly \(1\) at \(P\) in the rational-function
  picture: \(f \in K(C) \setminus \bar k\) with \(\mathrm{div}(f)_\infty = [P]\)).
  Therefore by \texttt{lem:morphism_to_p1_from_global_sections} there
  exists a unique morphism
  \[
    \varphi \;:=\; \varphi_{(1, f)} \colon C \longrightarrow \mathbb P^1_{\bar k}
  \]
  with \(\varphi^* x_0 = 1\) and \(\varphi^* x_1 = f\) (so the pullback of
  the standard ratio coordinate \(u = x_1 / x_0\) on $\mathbb P^1 \setminus
  \{\infty\}$ is $f$, and the pullback of the inverse coordinate
  \(t_\infty = x_0 / x_1 = 1 / f\) is the local parameter \(1/f\) at \(P\)).
  \emph{Step 3: \(\deg(\varphi) = 1\) via the pole divisor of \(f\).} The
  morphism \(\varphi\) is non-constant (since \(f\) is non-constant as a
  rational function on \(C\), the pullback \(\varphi^* u = f\) takes
  infinitely many values, ruling out a constant \(\varphi\)). By
  \texttt{lem:degree_via_pole_divisor} applied to \(\varphi\),
  \[
    \deg(\varphi)
    \;=\;
    \deg\!\left(\,\sum_{Q \in \varphi^{-1}(\infty)}
                  \ord_Q(\varphi^* t_\infty) \cdot [Q]\,\right).
  \]
  We compute the pole divisor of \(f\) as a rational function on \(C\):
  since \(f \in H^0(C, \mathcal O_C([P]))\) with \(f\) non-constant and
  \(\mathcal O_C([P])\) the line bundle of the effective divisor \([P]\), the
  rational-function \(f\) has pole divisor \((f)_\infty = [P]\) (the
  unique pole of order \(1\) at \(P\), no other poles; this is the standard
  identification of \(H^0(C, \mathcal O_C(D))\) with the rational functions
  \(g \in K(C)\) satisfying \(\mathrm{div}(g) + D \geq 0\), specialised to
  \(D = [P]\)). Hence \(\varphi^{-1}(\infty) = \{P\}\) with multiplicity~\(1\)
  and
  \[
    \deg(\varphi) \;=\; \deg([P]) \;=\; 1.
  \]
  \emph{Step 4: \(\varphi\) is an isomorphism.} By Hartshorne's IV.2
  identification \(\deg(\varphi) = [K(C) : \bar k(\mathbb P^1)]\), the
  degree-one conclusion of Step~3 rewrites as
  \(\operatorname{Module.finrank}_{\bar k(\mathbb P^1)} K(C) = 1\) (with
  the canonical \(\varphi\)-induced
  \texttt{[Algebra \,(ProjectiveLineBar kbar).functionField\,
   C.functionField]} instance). By
  \texttt{lem:degree_one_morphism_iso} applied to the non-constant
  morphism \(\varphi \colon C \to \mathbb P^1_{\bar k}\) (smooth proper
  geometrically irreducible curves over \(\bar k\) on both sides), under
  this function-field-extension-degree-one hypothesis, \(\varphi\) is an
  isomorphism of \(\bar k\)-schemes.
  This produces the required isomorphism \(C \cong \mathbb P^1_{\bar k}\)
  closing the headline target
  \texttt{AlgebraicGeometry.genusZero\_curve\_iso\_P1}.
  \emph{Comparison with Hartshorne's original argument.} Hartshorne's
  Example~IV.1.3.5 phrases the argument in terms of two distinct points
  \(P, Q \in C\) and the degree-zero divisor \(D = P - Q\), deriving \(P \sim Q\)
  and concluding via II.6.10.1 that \(C\) is rational; combined with I.6.12
  this yields \(C \cong \mathbb P^1\). The present proof restates the same
  combinatorial content via the more directly formalisable linear-system
  construction: instead of going through the \(P \sim Q\) statement and
  invoking II.6.10.1 as a black box, we produce the explicit morphism
  \(\varphi \colon C \to \mathbb P^1\) from the two sections \((1, f)\) of
  \(\mathcal O_C([P])\) given by Riemann--Roch and identify its degree
  directly. The two arguments are mathematically equivalent and have the
  same dependency footprint (Riemann--Roch in genus zero \(\Rightarrow\)
  \(\ell([P]) = 2\) \(\Rightarrow\) non-constant \(f\) \(\Rightarrow\) degree-\(1\)
  map to \(\mathbb P^1\) \(\Rightarrow\) iso); the linear-system phrasing maps
  more cleanly onto the Mathlib \texttt{Proj.fromOfGlobalSections} API and
  the project's existing Weil-divisor/degree pipeline of \texttt{RR.1}.
\end{proof}
\begin{remark}
[Drop of \texttt{sorryAx} on the genus-\(0\) Albanese keystone]
  \label{rmk:rr4_unblocks_sorryAx_on_rigidity_genus0}
  \uses{thm:genus_zero_curve_iso_p1}
  Closure of \texttt{chap:RiemannRoch_RationalCurveIso} retires the
  \texttt{sorryAx} propagation tagged on
  \texttt{AlgebraicGeometry.rigidity\_genus0\_curve\_to\_grpScheme}
  (\texttt{thm:rigidity_genus0_curve_to_AV} of
  \texttt{chap:avr_for_rr}, body landed iter-166), whose
  \emph{first} of two upstream \texttt{sorry} sources is the
  \texttt{genusZero\_curve\_iso\_P1} target this chapter closes. The
  \emph{second} upstream is the helper
  \texttt{morphism\_P1\_to\_grpScheme\_const}'s residual sorries, tracked
  separately under the genus-\(0\) Albanese arm of \texttt{Jacobian.tex}.
\end{remark}
\section{Out of scope}
The following items are deliberately \emph{not} part of \texttt{RR.4} and
are sequenced elsewhere:
\begin{itemize}
  \item The general ``smooth proper geometrically irreducible curve \(C\) of
    genus \(0\) over a \emph{non-closed} field \(k\) with a \(k\)-rational point
    \(\cong \mathbb P^1_k\)'' statement is the \(\bar k \to k\) descent of
    \texttt{chap:RiemannRoch_RationalCurveIso}. The descent is hosted in
    \texttt{Jacobian.tex} via \texttt{genusZeroWitness.key}; it consumes
    \texttt{chap:RiemannRoch_RationalCurveIso} together with Galois descent for
    smooth proper curves and is out of scope for \texttt{RR.4}.
  \item The Mathlib API \texttt{AlgebraicGeometry.Proj.fromOfGlobalSections}
    (\texttt{Mathlib.AlgebraicGeometry.ProjectiveSpectrum.Basic}) is
    consumed as a black box by
    \texttt{lem:morphism_to_p1_from_global_sections}; its own proof inside
    Mathlib (the universal property of \(\mathrm{Proj}\), graded-ring
    presentation of \(\mathbb P^n\)) is not re-derived here.
  \item Multiplicativity of degree under finite pullback (Hartshorne~II.6.9)
    is consumed as a black box by \texttt{lem:degree_via_pole_divisor};
    its proof is a separable sub-lemma threaded inside
    \texttt{AlgebraicJacobian/RiemannRoch/RationalCurveIso.lean} and is
    \emph{not} part of \texttt{RR.4}'s blueprint statement (it lives at
    the same logical level as the closely related
    \texttt{thm:principal_deg_zero} of \texttt{RR.1}).
  \item The full equivalence of categories ``smooth projective curves over
    \(\bar k\) with dominant morphisms'' \(\cong\) ``function fields of
    transcendence degree~\(1\) over \(\bar k\) with \(\bar k\)-homomorphisms''
    (Hartshorne~I.6.12) is consumed as a high-level framing in
    \texttt{lem:degree_one_morphism_iso} but is \emph{not} re-formalised at
    the equivalence-of-categories level here. The proof body in the Lean
    file threads the spelled-out scheme-theoretic ``finite-degree-one
    morphism between integral normal schemes is an isomorphism''
    argument (Stacks~\textbf{0AVX}), which is the minimal-Mathlib-dependence
    pin compatible with the chapter's downstream needs.
  \item ``A smooth proper geometrically irreducible curve over \(\bar k\)
    has at least two distinct \(\bar k\)-points'' (the non-empty / infinite
    cardinality of \(C(\bar k)\)) is consumed implicitly by the present
    chapter but its formal statement lives in the project's curve API
    (\texttt{Genus.tex} / \texttt{AlgebraicJacobian.Genus}); the chapter
    does not restate it.
  \item Cohomology of the line bundle \(\mathcal O_C([P])\) on a curve
    (Hartshorne~III.5 + IV.1) is fully consumed via \texttt{RR.3}; in
    particular the \(H^1\)-vanishing of \(\mathcal O_C([P])\) on a genus-\(0\)
    curve is part of \texttt{RR.3}'s payload feeding
    \texttt{thm:riemannRoch_genus_zero} of \texttt{RR.2}, and is not
    re-derived here.
\end{itemize}
